On simulated noisy Lorenz-63 with a fixed library and STLSQ, the registered claim that the weak form beats every derivative-based method in exact-support rate at every noise level from 1\% is refuted: it ties three baselines at 1\%, its leads at 2\% and 5\% are significant against some baselines only, and at 10\% every method fails. The weak form has the lowest mean coefficient error at every level, significantly against all four baselines only at 0\% and 10\%; the clearest separation, at 5\%, is in the median (0.0129 against 0.0456--0.0739), while the mean there is not significantly below FD or TV. At the tuned thresholds no smoother lowered the mean coefficient error of finite differences at 5\%, which refutes the second claim as registered; that verdict depends on the finite-difference threshold, and at equal thresholds only Savitzky--Golay is above finite differences throughout while TV is below. Up to 1\% noise the support rate is governed by the STLSQ threshold more than by the derivative source, and a tie rule in the tuner was enough to reverse the low-noise ordering; at 5\% the weak form's rate is above the best threshold of each baseline. These conclusions hold for one simulated system, 20 trials per cell and oracle-assisted tuning.
